\documentclass[10pt,a4paper]{amsart}
\usepackage{amsmath,amssymb,amsthm,graphicx}
% Retained minimal preamble for the Stacks Project chapter more-algebra.tex (More on Algebra).
% The chapter file pinned-original-source.tex carries only \input{preamble} on its first line, so no
% preamble text lives inside the chapter itself.  The Stacks Project uses one single shared preamble for
% every chapter; the snapshot stages it identically as sources/stacks-chow/preamble.tex and
% sources/stacks-etale-cohomology/preamble.tex.  The block below is copied mechanically, line for line,
% from sources/stacks-chow/preamble.tex lines 14-18 (the Xy-pic setup) and 154-217 (the native theorem
% declarations together with the mathematical macros).  Nothing was added, removed or redefined; the
% wrapper class and the amsmath/amssymb/amsthm/graphicx base packages are supplied by the assembler.
% Chapter: more-algebra.tex -- More on Algebra.
% For commutative diagrams we use Xy-pic
\usepackage[all]{xy}

% We use 2cell for 2-commutative diagrams.
\xyoption{2cell}

% Theorem environments.
%
\theoremstyle{plain}
\newtheorem{theorem}[subsection]{Theorem}
\newtheorem{proposition}[subsection]{Proposition}
\newtheorem{lemma}[subsection]{Lemma}

\theoremstyle{definition}
\newtheorem{definition}[subsection]{Definition}
\newtheorem{example}[subsection]{Example}
\newtheorem{exercise}[subsection]{Exercise}
\newtheorem{situation}[subsection]{Situation}

\theoremstyle{remark}
\newtheorem{remark}[subsection]{Remark}
\newtheorem{remarks}[subsection]{Remarks}

\numberwithin{equation}{subsection}

% Macros
%
\def\lim{\mathop{\mathrm{lim}}\nolimits}
\def\colim{\mathop{\mathrm{colim}}\nolimits}
\def\Spec{\mathop{\mathrm{Spec}}}
\def\Hom{\mathop{\mathrm{Hom}}\nolimits}
\def\Ext{\mathop{\mathrm{Ext}}\nolimits}
\def\SheafHom{\mathop{\mathcal{H}\!\mathit{om}}\nolimits}
\def\SheafExt{\mathop{\mathcal{E}\!\mathit{xt}}\nolimits}
\def\Sch{\mathit{Sch}}
\def\Mor{\mathop{\mathrm{Mor}}\nolimits}
\def\Ob{\mathop{\mathrm{Ob}}\nolimits}
\def\Sh{\mathop{\mathit{Sh}}\nolimits}
\def\NL{\mathop{N\!L}\nolimits}
\def\CH{\mathop{\mathrm{CH}}\nolimits}
\def\proetale{{pro\text{-}\acute{e}tale}}
\def\etale{{\acute{e}tale}}
\def\QCoh{\mathit{QCoh}}
\def\Ker{\mathop{\mathrm{Ker}}}
\def\Im{\mathop{\mathrm{Im}}}
\def\Coker{\mathop{\mathrm{Coker}}}
\def\Coim{\mathop{\mathrm{Coim}}}

% Boxtimes
%
\DeclareMathSymbol{\boxtimes}{\mathbin}{AMSa}{"02}

%
% Macros for moduli stacks/spaces
%
\def\QCohstack{\mathcal{QC}\!\mathit{oh}}
\def\Cohstack{\mathcal{C}\!\mathit{oh}}
\def\Spacesstack{\mathcal{S}\!\mathit{paces}}
\def\Quotfunctor{\mathrm{Quot}}
\def\Hilbfunctor{\mathrm{Hilb}}
\def\Curvesstack{\mathcal{C}\!\mathit{urves}}
\def\Polarizedstack{\mathcal{P}\!\mathit{olarized}}
\def\Complexesstack{\mathcal{C}\!\mathit{omplexes}}
% \Pic is the operator that assigns to X its picard group, usage \Pic(X)
% \Picardstack_{X/B} denotes the Picard stack of X over B
% \Picardfunctor_{X/B} denotes the Picard functor of X over B
\def\Pic{\mathop{\mathrm{Pic}}\nolimits}
\def\Picardstack{\mathcal{P}\!\mathit{ic}}
\def\Picardfunctor{\mathrm{Pic}}
\def\Deformationcategory{\mathcal{D}\!\mathit{ef}}
\usepackage{cleveref}
\crefname{originalproposition}{original proposition}{original propositions}
\Crefname{originalproposition}{Original proposition}{Original propositions}
\makeatletter
\newlabel{algebra-proposition-unramified-locally-standard}{{0395}{source}{}{sourcepointer-1}{}}
\newlabel{algebra-proposition-unramified-locally-standard@cref}{{[originalproposition][0][]0395}{[1][0][]source}{}{}{}}
\makeatother
\crefname{originallemma}{original lemma}{original lemmas}
\Crefname{originallemma}{Original lemma}{Original lemmas}
\makeatletter
\newlabel{algebra-lemma-etale-over-field}{{00U3}{source}{}{sourcepointer-2}{}}
\newlabel{algebra-lemma-etale-over-field@cref}{{[originallemma][0][]00U3}{[1][0][]source}{}{}{}}
\makeatother
\crefname{originallemma}{original lemma}{original lemmas}
\Crefname{originallemma}{Original lemma}{Original lemmas}
\makeatletter
\newlabel{lemma-invariance-number-branches-smooth}{{0DQ1}{source}{}{sourcepointer-3}{}}
\newlabel{lemma-invariance-number-branches-smooth@cref}{{[originallemma][0][]0DQ1}{[1][0][]source}{}{}{}}
\makeatother
\crefname{originallemma}{original lemma}{original lemmas}
\Crefname{originallemma}{Original lemma}{Original lemmas}
\makeatletter
\newlabel{lemma-number-of-branches-1}{{0C37}{source}{}{sourcepointer-4}{}}
\newlabel{lemma-number-of-branches-1@cref}{{[originallemma][0][]0C37}{[1][0][]source}{}{}{}}
\makeatother
\crefname{originallemma}{original lemma}{original lemmas}
\Crefname{originallemma}{Original lemma}{Original lemmas}
\makeatletter
\newlabel{algebra-lemma-silly}{{00DS}{source}{}{sourcepointer-5}{}}
\newlabel{algebra-lemma-silly@cref}{{[originallemma][0][]00DS}{[1][0][]source}{}{}{}}
\makeatother
\crefname{originallemma}{original lemma}{original lemmas}
\Crefname{originallemma}{Original lemma}{Original lemmas}
\makeatletter
\newlabel{lemma-ideal-inverse-nonzero}{{0GS5}{source}{}{sourcepointer-6}{}}
\newlabel{lemma-ideal-inverse-nonzero@cref}{{[originallemma][0][]0GS5}{[1][0][]source}{}{}{}}
\makeatother
\crefname{originallemma}{original lemma}{original lemmas}
\Crefname{originallemma}{Original lemma}{Original lemmas}
\makeatletter
\newlabel{lemma-reduced-goes-up}{{07QK}{source}{}{sourcepointer-7}{}}
\newlabel{lemma-reduced-goes-up@cref}{{[originallemma][0][]07QK}{[1][0][]source}{}{}{}}
\makeatother
\title{A map that is injective and unramified at a prime is etale after a principal localization, over a geometrically unibranch domain}
\author{The Stacks Project Authors}
\date{stacks-project:more-algebra.tex:local-snapshot-20261012; GNU FDL 1.2 or later}
\begin{document}
\maketitle
\setcounter{subsection}{1}
\noindent\textit{Editorial note.} All delivered mathematical text is the author's own, copied from the pinned Stacks Project chapter file more-algebra.tex; no statement or proof step was rewritten, reordered or removed. The counted statement is the labeled lemma environment at source lines 31762-31774 (canonical tag 0GSC) and the counted proof is the proof environment at source lines 31776-31815, which immediately follows it and is its own proof.\par\medskip
\noindent\textit{Editorial note.} The proof refers to four results of other chapters, namely the locally standard form of unramified maps, etale algebras over a field, the element separating minimal primes and normal algebras for etale maps, and to three results of this chapter among which are the two branch-count statements and the lemma on ideals meeting the source domain. All are carried as navigation stubs with their roles recorded in the reasons.\par\medskip
\noindent\textit{Editorial note.} The proof contains the slip that the minimal primes of the etale cover are lie over the generic point of the base, which is copied exactly as printed and is not corrected.\par\medskip
\noindent\textit{Editorial note.} The chapter file pinned-original-source.tex carries only an input of the shared preamble on its first line, so no preamble text lives inside the chapter. The retained file sources/stacks-more-algebra/retained-preamble.tex is a mechanical, line-for-line copy of the project's single shared preamble as staged in the snapshot (Xy-pic setup plus the native theorem declarations and mathematical macros). Nothing was added, removed or redefined, and the excerpt serves only as preamble, never as a span coordinate.\par\medskip
\noindent\textit{Editorial note.} The source URL is the source-level Stacks Project URL, as required when more than one card is produced from one source; the per-claim tag is carried by claim\_key and by this note. The counted statement and proof cite no bibliography entry, so this card needs no bibliography.\par\medskip
\noindent\textit{Allowed standard prerequisites (editorial).}\begin{itemize}
\item An unramified ring map of finite type is locally standard etale, so after a principal localization there is an etale map onto the ring with a surjection from it. invoked at proof lines 31777-31779. Replaces the given map by a surjection from a standard etale algebra, which is the setting in which the injectivity argument is run.
\item An etale algebra over a field is a finite product of fields, and an etale map is flat with the same minimal primes as the base, so the number of minimal primes is finite. used at proof lines 31781-31788. Identifies the fibre over the generic point and shows that the etale cover has finitely many minimal primes.
\item Over a geometrically unibranch local ring the localization of an etale algebra at a prime is a domain, so the prime contains exactly one minimal prime. used at proof lines 31788-31797. Justifies the choice of an element that vanishes on all but one minimal prime so that the cover becomes a domain after localization.
\end{itemize}\medskip
\ifdefined\hypertarget\hypertarget{sourcepointer-1}{}\fi
\noindent\textit{Source reference.} Original proposition 0395: sources/stacks-chow/tags.txt, source lines 4260--4260. The reference produces the standard etale cover with a surjection onto the given ring, used at proof lines 31777-31779; the reduction and the injectivity argument that follow are inside the counted span.\par\medskip
\ifdefined\hypertarget\hypertarget{sourcepointer-2}{}\fi
\noindent\textit{Source reference.} Original lemma 00U3: sources/stacks-chow/tags.txt, source lines 1097--1097. The reference identifies the fibre of an etale algebra over a field as a finite product of fields, used at proof lines 31784-31786 and again at the end for the residue fields; the reduction built on it is inside the counted span.\par\medskip
\ifdefined\hypertarget\hypertarget{sourcepointer-3}{}\fi
\noindent\textit{Source reference.} Original lemma 0DQ1: sources/stacks-more-algebra/pinned-original-source.tex, source lines 31596--31596. The reference is one of the two applied at proof lines 31790-31791 to see that the localization of the etale cover at the prime is a domain; the use of that domain property is inside the counted span.\par\medskip
\ifdefined\hypertarget\hypertarget{sourcepointer-4}{}\fi
\noindent\textit{Source reference.} Original lemma 0C37: sources/stacks-more-algebra/pinned-original-source.tex, source lines 31566--31566. The reference is the second of the two applied at proof lines 31790-31791 for the same purpose; it is a pointer to a same-chapter result and everything drawn from it is inside the counted span.\par\medskip
\ifdefined\hypertarget\hypertarget{sourcepointer-5}{}\fi
\noindent\textit{Source reference.} Original lemma 00DS: sources/stacks-chow/tags.txt, source lines 508--508. The reference supplies the element that lies in all minimal primes except one, used at proof lines 31794-31795; the localization that consumes it is inside the counted span.\par\medskip
\ifdefined\hypertarget\hypertarget{sourcepointer-6}{}\fi
\noindent\textit{Source reference.} Original lemma 0GS5: sources/stacks-more-algebra/pinned-original-source.tex, source lines 31702--31702. The reference is applied at proof lines 31807-31809 to the domain map of localizations to find a nonzero element of the base mapping to zero, which is the contradiction; the setup of that contradiction is inside the counted span. The lemma itself is the subject of a separate card of this lane with disjoint coordinates.\par\medskip
\ifdefined\hypertarget\hypertarget{sourcepointer-7}{}\fi
\noindent\textit{Source reference.} Original lemma 07QK: sources/stacks-more-algebra/pinned-original-source.tex, source lines 10533--10533. The reference is cited at proof lines 31783-31784 to see that the etale cover of the domain is reduced, which the reduction to a domain then uses; the reduction itself is inside the counted span. The author cites it with the label as printed, which lacks the chapter prefix of the other Algebra references in the same proof.\par\medskip
\begin{lemma}
\label{lemma-local-unramified-extension-unibranch-domain-is-etale}
Let $A \to B$ be a ring map. Let $\mathfrak q \subset B$ be a prime
ideal lying over the prime $\mathfrak p \subset A$. Assume
\begin{enumerate}
\item $A$ is a domain,
\item $A_\mathfrak p$ is geometrically unibranch,
\item $A \to B$ is unramified at $\mathfrak q$, and
\item $A_\mathfrak p \to B_\mathfrak q$ is injective.
\end{enumerate}
Then there exists a $g \in B$, $g \not \in \mathfrak q$ such that
$B_g$ is \'etale over $A$.
\end{lemma}

\begin{proof}
By Algebra, Proposition \ref{algebra-proposition-unramified-locally-standard}
after replacing $B$ by a principal localization, we can find a
standard \'etale ring map $A \to B'$ and a surjection $B' \to B$.
Denote $\mathfrak q' \subset B'$ the inverse image of $\mathfrak q$.
We will show that $B' \to B$ is injective
after possibly replacing $B'$ by a principal localization.

\medskip\noindent
In this paragraph we reduce to the case that $B'$ is a domain.
Since $A$ is a domain, the ring $B'$ is reduced, see
Algebra, Lemma \ref{lemma-reduced-goes-up}.
Let $K$ be the fraction field of $A$. Then $B' \otimes_A K$
is \'etale over a field, hence is a finite product of fields, see
Algebra, Lemma \ref{algebra-lemma-etale-over-field}.
Since $A \to B'$ is \'etale (hence flat) the minimal primes of $B'$
are lie over $(0) \subset A$ (by going down for flat ring maps).
We conclude that $B'$ has finitely many minimal primes, say
$\mathfrak r_1, \ldots, \mathfrak r_r \subset B'$.
Since $A_\mathfrak p$ is geometrically unibranch and $A \to B'$ \'etale,
the ring $B'_{\mathfrak q'}$ is a domain, see
Lemmas \ref{lemma-invariance-number-branches-smooth} and
\ref{lemma-number-of-branches-1}.
Hence $\mathfrak q' \supset \mathfrak r_i$ for exactly one $i = i_0$.
Choose $g' \in B'$, $g' \not \in \mathfrak r_{i_0}$ but
$g' \in \mathfrak r_i$ for $i \not = i_0$, see
Algebra, Lemma \ref{algebra-lemma-silly}. After replacing
$B'$ and $B$ by $B'_{g'}$ and $B_{g'}$ we obtain that $B'$
is a domain.

\medskip\noindent
Assume $B'$ is a domain, in particular $B' \subset B'_{\mathfrak q'}$.
If $B' \to B$ is not injective,
then $J = \Ker(B'_{\mathfrak q'} \to B_\mathfrak q)$ is nonzero.
By Lemma \ref{lemma-ideal-inverse-nonzero} applied to
$A_\mathfrak p \to B'_{\mathfrak q'}$
we find a nonzero element $a \in A_\mathfrak p$
mapping to zero in $B_\mathfrak q$ contradicting
assumption (4). This finishes the proof.
\end{proof}

\end{document}
